\documentclass[10pt,a4paper]{article}
\usepackage[OT1]{fontenc}
\usepackage[utf8]{inputenc}
\usepackage[ngerman]{babel}
\usepackage[a4paper,margin=13mm,top=12mm,bottom=14mm]{geometry}
\usepackage{microtype}
\usepackage{enumitem}
\usepackage{titlesec}
\usepackage{fancyhdr}
\usepackage[hidelinks]{hyperref}
\setlength{\parindent}{0pt}
\setlength{\parskip}{1pt}
\setlist[itemize]{label=$\bullet$,leftmargin=4mm,itemsep=0.6pt,parsep=0pt,topsep=1.5pt}
\titleformat{\section}{\large\bfseries}{}{0pt}{}
\titlespacing*{\section}{0pt}{4pt}{4pt}
\newcommand{\block}[2]{\par\vspace{3pt}{\bfseries #1}\hfill{\scriptsize #2}\par\nopagebreak}
\newcommand{\qa}[3]{\par\vspace{2.5pt}\textbf{#1. #2}\par\nopagebreak #3\par}
\pagestyle{fancy}
\fancyhf{}
\fancyhead[L]{\small Security Engineering -- Lecture 1 (16.10.2026) -- Minitest 23.10.2026}
\fancyhead[R]{\small Lernzettel}
\fancyfoot[L]{\scriptsize F = PDF-Seite Lecture1; P = PDF-Seite Preamble (Lecture 0); T = Transkript-Zeit; M = Moodle-Kursseite; \textbf{?F77} = Frage von dieser Fragenfolie}
\fancyfoot[R]{\small\thepage}
\setlength{\headheight}{13pt}
\begin{document}
\section*{1\quad Muss ich wissen -- Grundlagen}
\block{Security Engineering: Definition und Abgrenzung}{F5--7; T1:41--4:53 \quad \textbf{?F77}}
\begin{itemize}
\item \textbf{Security Engineering} baut Systeme, die trotz \textbf{malice, error and mischance} dependable bleiben, mit Tools, Prozessen und Methoden für design, implement, test und Anpassung an eine veränderte Umgebung.
\item \textbf{Security vs.\ Safety}: Bei Security ist meist ein \textbf{Threat Actor} beteiligt, ohne Angreifer geht es um Safety.
\item \textbf{Software Engineering} sorgt dafür, dass bestimmte Dinge passieren (\glqq John can read this file\grqq{}), Security ist aber nicht bloß das Verhindern von Dingen, weil die Realität komplexer ist.
\item \textbf{Security Requirements} sind systemabhängig und meist eine Kombination aus user authentication, transaction integrity, accountability, fault-tolerance und message secrecy.
\item \textbf{Adversarial Thinking} heißt, frühere Angriffe, Entry Paths, Entwicklungen und Outcomes zu studieren.
\end{itemize}
\block{Disziplinen und warum Systeme scheitern}{F6,8,16; T2:21--3:07,5:04,6:14 \quad \textbf{?F79}}
\begin{itemize}
\item \textbf{Cross-disciplinary expertise}: Business Processes, Software Engineering und Testing, Cryptography und Computer Security bis Hardware Tamper-Resistance sowie Economics, Applied Psychology, Organizations und Law.
\item \textbf{Hauptgrund für Fehlschläge}: Designer schützen die falschen Dinge oder die richtigen Dinge auf die falsche Weise.
\item \textbf{Zu enge Systemgrenze}: Als sicher beworbene Systeme brachen, sobald eine bestimmte Application lief oder das Gerät anders als vorgesehen genutzt wurde.
\item \textbf{Human Components ignorieren} und damit Usability vernachlässigen ist laut F16 eine der größten Ursachen für Security Failure.
\item \textbf{Falsche Incentives}: Wer das System bewacht und wartet, muss geschult und richtig motiviert sein, sonst werden die Schutzziele verfehlt.
\end{itemize}
\block{System: sechs Definitionen -- im Kurs gilt Definition 6}{F16; T13:08--15:13 \quad \textbf{?F77}}
\begin{itemize}
\item \textbf{1} = Produkt oder Komponente wie ein kryptografisches Protokoll, eine Smartcard oder die Hardware von Phone, Laptop oder Server.
\item \textbf{2} = 1 plus Operating System, Communications und weitere Infrastruktur; \textbf{3} = plus Applications (Banking-App, Browser, LLM \dots{}, Client- und Cloud-Teile).
\item \textbf{4} = plus IT Staff; \textbf{5} = plus interne Users und Management; \textbf{6} = plus Customers und andere externe Users.
\item \textbf{Merksatz Vorlesung}: Krypto-Leute denken meist an 1--2, Firmen hören bei 5 auf, wir nehmen alle sechs Stufen.
\end{itemize}
\block{Workflow: vier Schritte in dieser Reihenfolge}{F10--13,32; T7:19--12:58 \quad \textbf{?F77}}
\begin{itemize}
\item \textbf{1.\ Threat Model}: Wer könnte warum und wie angreifen -- Insider oder Outsider, People oder Malware, Government, Crooks oder Script Kiddy?
\item \textbf{2.\ Security Policy}: Welche Protection Properties sollen gelten (Keep secret? Guard money?) und welche Regeln brauchen wir?
\item \textbf{3.\ Implementation}: Welche Protection Mechanisms setzen wir ein und warum?
\item \textbf{4.\ Assurance}: Woher wissen wir, dass wir genug getan haben, und wie beweisen wir das Dritten?
\item \textbf{Assurance $\neq$ Compliance}: Compliance ist nur das Zertifikat (z.\,B. Common Criteria, ISO 27000), Assurance stellt die Zielerreichung wirklich sicher und weist sie Dritten nach.
\end{itemize}
\block{Threat Model: warum wichtig?}{F33; T7:46--8:55,40:31--41:10 \quad \textbf{?F77}}
\begin{itemize}
\item \textbf{Likely Threat Actors} zuerst identifizieren, jeweils mit ihren \textbf{Incentives/Goals} und \textbf{Capabilities/Knowledge}.
\item \textbf{Kein Totalschutz}: Gegen alle denkbaren Threats schützen und trotzdem nützliche Arbeit zu reasonable cost leisten geht nicht, also priorisiert man die dringendsten Threat Actors.
\item \textbf{Kontext entscheidet}: Ein Chiphersteller wie Infineon fürchtet eher Nation-State Actors und Konkurrenz, Payment/Banking eher Cyberkriminelle mit monetären Incentives.
\item \textbf{Insider Threats} sind Menschen, deren Incentives nicht mehr zur Firma passen, und laut Vorlesung auch Malware.
\end{itemize}
\block{Security Engineering Analysis Framework (Figure 1.1)}{F9,14; T6:55--7:07,12:19--12:58}
\begin{itemize}
\item \textbf{Policy} = was das System erreichen soll, \textbf{Mechanism} = die technischen Mittel wie Ciphers und Cryptography.
\item \textbf{Assurance} = wie viel Verlass auf die Mechanismen besteht, \textbf{Incentives} = die Motive derer, die das System bewachen oder angreifen.
\item \textbf{Alle vier Kästen sind miteinander verbunden} und beeinflussen sich gegenseitig.
\end{itemize}

\section*{1\quad Muss ich wissen -- Definitionen}
\block{Person, Subject, Principal}{F17--18; T15:18--18:14 \quad \textbf{?F78}}
\begin{itemize}
\item \textbf{Person} = physical person oder legal person wie eine Firma oder eine Regierung.
\item \textbf{Subject} = physical person in irgendeiner Rolle, z.\,B.\ Operator, Principal oder Victim, also nie eine Firma.
\item \textbf{Principal} = Entität, die authentifiziert werden kann und in einem Security System handelt, z.\,B.\ Subject, Person, Role oder Equipment wie Laptop, Phone, Smartcard oder Communication Channel (Port-Nr.).
\item \textbf{Zusammengesetzte Principals}: Group (Alice or Bob), Conjunction (Alice and Bob acting together), Compound Role (Alice acting as Bob's manager) und Delegation (Bob acting for Alice in her absence).
\item \textbf{Alice-und-Bob-Konvention} ist ungenau: Wer ist bei \glqq Alice authenticates herself to Bob\grqq{} gemeint, Mensch, Key oder Software?
\end{itemize}
\block{Group vs.\ Role und Levels of Abstraction}{F19; T18:14--19:44 \quad \textbf{?F78}}
\begin{itemize}
\item \textbf{Group} = Menge von Principals, etwa alle Studierenden im Hörsaal.
\item \textbf{Role} = Menge von Funktionen, die nacheinander von verschiedenen Personen übernommen werden, etwa \glqq the rector of AAU\grqq{} oder \glqq student\grqq.
\item \textbf{Levels of Abstraction}: \glqq Bob's smartcard representing Bob who is acting for Alice\grqq{} ist etwas anderes als \glqq Bob operating Alice's smartcard in her absence\grqq.
\item \textbf{Regel}: präzise beschreiben und das Abstraction Level ausdrücklich angeben.
\end{itemize}
\block{Identity, Trust, Trustworthiness}{F20--21; T19:44--22:08 \quad \textbf{?F77}}
\begin{itemize}
\item \textbf{Identity} = Correspondence zwischen den Namen zweier Principals, die anzeigt, dass beide dieselbe Person oder dasselbe Equipment meinen, also kein bloßer Name.
\item \textbf{Beispiel}: Bob in \glqq Bob acting for Alice in her absence\grqq{} ist derselbe Bob wie in \glqq Bob acting as Charlie's manager\grqq.
\item \textbf{Trusted} = System oder Komponente, deren Versagen die Security Policy brechen kann, z.\,B.\ das TPM im Computer.
\item \textbf{Trustworthy} = System oder Komponente, die nicht versagt.
\item \textbf{Trusted but not trustworthy}: ein NSA-Mitarbeiter, der unautorisiert Key Material an einen chinesischen Diplomaten verkauft.
\end{itemize}

\block{Confidentiality, Privacy, Anonymity}{F22--23; T22:29--26:08 \quad \textbf{?F78}}
\begin{itemize}
\item \textbf{Confidentiality (non-disclosure)} = Verpflichtung, Geheimnisse einer anderen Person oder Organisation zu schützen, wenn man sie kennt.
\item \textbf{Privacy} = Fähigkeit und/oder Recht, die eigenen persönlichen Informationen zu schützen, erweitert auf das Verhindern von Eingriffen in den persönlichen Raum.
\item \textbf{Merksatz (Secrecy = Preventing Information Sharing)}: Privacy ist Secrecy zum Nutzen des Individuums, Confidentiality zum Nutzen der Organisation (Patient = Privacy, Pflicht des Arztes = Confidentiality).
\item \textbf{Anonymity} = lack of attribution einer Person zu einer von ihr begangenen Tat, bis hin zu bewusster Geheimhaltung.
\item \textbf{Internet nicht für Anonymity entworfen}: Protokolle und IP-Adressen verraten Metadaten, die leichter zu bekommen sind als Inhalte, und anonym ist man nur in einer Crowd to hide in.
\end{itemize}
\block{Integrity, Authenticity, Availability, Security Failure}{F24; T26:42--28:09 \quad \textbf{?F78}}
\begin{itemize}
\item \textbf{Integrity} = Zusicherung, dass Daten nicht verändert wurden, das Gegenteil von Data Corruption (at move/rest).
\item \textbf{Authenticity} = Integrity \textbf{plus Freshness plus Proof of Identity}: Man spricht nachweislich mit einem genuine Principal und nicht mit einem Replay alter Nachrichten.
\item \textbf{Freshness-Beispiel}: Zertifikate für digitale Signaturen gelten nur eine bestimmte Zeit.
\item \textbf{Availability} = Daten und Netzwerkressourcen sind für autorisierte Nutzer bei Bedarf zugänglich.
\item \textbf{Security Failure} = unerwünschte Verletzung der Security Policy eines Systems.
\end{itemize}
\block{Weakness, Vulnerability, Exploit, CWE/CVE}{F25--26; T29:03--32:51 \quad \textbf{?F78}}
\begin{itemize}
\item \textbf{Weakness} = Bedingung in einer SW/HW-Komponente, die zur Entstehung von Vulnerabilities beitragen könnte.
\item \textbf{Vulnerability} = Schwäche eines Systems oder seiner Umgebung, die zusammen mit einer internen oder externen Threat zu einem Security Failure führen kann.
\item \textbf{Exploit} = Methode oder Code, der eine Vulnerability ausnutzt.
\item \textbf{CWE} listet Weakness-Typen, z.\,B.\ \textbf{CWE-125 Out-of-bounds Read}: Das Produkt liest hinter dem Ende oder vor dem Anfang des vorgesehenen Buffers.
\item \textbf{CVE} listet konkrete Vulnerabilities, z.\,B.\ \textbf{CVE-2014-0160 Heartbleed} in der OpenSSL Library: improper input validation (missing bounds check) in der TLS Heartbeat Extension.
\end{itemize}
\block{Threat, Threat Actor, Hack}{F25,27; T28:24--29:03,32:51--33:45 \quad \textbf{?F79}}
\begin{itemize}
\item \textbf{Threat} = Handlung oder Ereignis, das zu einer unerwünschten negativen Wirkung auf ein Computersystem führen würde (Security Failure).
\item \textbf{Threat Actor} = Person, Gruppe, Organisation, Staat oder andere Entität mit der Fähigkeit, eine Threat zu cause, carry, transmit, support oder exploit.
\item \textbf{Zusammenspiel}: Ein Threat Agent nutzt per Exploit eine Vulnerability aus, dazwischen liegen Countermeasures.
\item \textbf{Hack nach Bruce Schneier} = etwas, das die Regeln eines Systems erlauben, das aber von seinen Designern weder vorhergesehen noch gewollt war.
\end{itemize}
\block{Security Policy, Security Target, Protection Profile, Protection}{F29--30; T35:53--39:32 \quad \textbf{?F79}}
\begin{itemize}
\item \textbf{Security Policy (what)} = knappe Aussage der Schutzstrategie eines Systems in Business Goals, z.\,B.\ \glqq in jeder Transaction sind Credits und Debits gleich, und alles über $\$$1.000.000 müssen zwei Manager autorisieren\grqq.
\item \textbf{Security Target/Goal (how)} = detailliertere Spezifikation (keine Implementation), wie die Policy in einem bestimmten Produkt umgesetzt wird (Encryption, Digital Signatures, Access Controls, Audit Logs), und Maßstab für die Implementation.
\item \textbf{Protection Profile} = Interface zwischen Security Target und Security Policy: wie ein Security Goal, aber device-independent genug für vergleichende Evaluations verschiedener Produkte und Versionen.
\item \textbf{Protection} = abstrakte Eigenschaft wie Confidentiality oder Integrity, abstrakt genug für allgemeine Systeme statt konkreter Implementierungen.
\item \textbf{Hierarchie} (abstrakt $\rightarrow$ konkret): Protection $\rightarrow$ Protection Profile $\rightarrow$ Security Target $\rightarrow$ Implementation.
\end{itemize}

\section*{1\quad Muss ich wissen -- Adversaries}
\block{Wer sind die Gegner und was treibt sie an?}{F34,61--62,67,75; T9:08--9:28,41:41--42:37 \quad \textbf{?F80 ?F82}}
\begin{itemize}
\item \textbf{State Actors} (Five Eyes, China, Russia, Iran, North Korea \dots) wollen den \textbf{Intelligence Contest} gewinnen, also Informations- und Machtvorteile.
\item \textbf{Crooks} sind Kriminelle mit \textbf{economic incentives}, z.\,B.\ Ransomware-Gangs und Fraudsters, von kleinem Betrug bis zu großen Cryptocurrency Hacks.
\item \textbf{The Swamp} sind Kriminelle mit \textbf{abuse incentives}, z.\,B.\ Hate Crimes, Abuse und Bullying.
\item \textbf{Lawful Operators} wie Security Researchers und Tool Vendors kommen nächstes Mal, \textbf{Hacktivism und Script Kiddies} werden nicht behandelt.
\item \textbf{\glqq The rest of us\grqq}: Wir alle können in verschiedenen Rollen unwissentlich als Amplifier für schädliches Material wirken.
\end{itemize}
\block{Intelligence Contests (Intelligence Doctrine)}{F61; T1:23:10--1:23:50 \quad \textbf{?F81}}
\begin{itemize}
\item \textbf{Quelle}: Joshua Rovner, früher scholar-in-residence der NSA.
\item \textbf{1.\ Collect, 2.\ Exploit}: Rennen, mehr und bessere Informationen als der Gegner zu sammeln und damit die eigene Position zu verbessern.
\item \textbf{3.\ Undermine}: wechselseitiger Versuch, Moral, Institutionen und Allianzen des Gegners zu untergraben.
\item \textbf{4.\ Disable, 5.\ Preposition}: gegnerische Fähigkeiten durch Sabotage ausschalten und Assets für den Konfliktfall vorpositionieren.
\item \textbf{Merksatz}: Intelligence Contests werden nie wirklich gewonnen oder verloren.
\end{itemize}
\block{Five Eyes (FVEY): wer und wie?}{F36--37,45; T43:05--46:08,56:53--58:19,1:00:50 \quad \textbf{?F80}}
\begin{itemize}
\item \textbf{Mitglieder}: USA, UK, Canada, Australia und New Zealand teilen Intelligence-Infrastruktur.
\item \textbf{Platform Advantage}: Viele Glasfasern folgen alten Telegraph-/Telefonkabeln (Commonwealth-Routen), daher bulk wiretap u.\,a.\ in Cornwall, Gibraltar, Malta und Singapore.
\item \textbf{Weitere Quellen}: Satellite Downlinks und Embassies, und SIGINT-Agencies (NSA, GCHQ) bekommen mehr Geld als klassisches Human Spying.
\item \textbf{Crypto Wars}: Die Agencies wollten starke Kryptografie untergraben, z.\,B.\ mit dodgy RNGs und sabotierten Protokollstandards, um VPN-Traffic zu entschlüsseln.
\item \textbf{2013 Edward Snowden} enthüllte das Ausmaß, laut Vorlesung als Whistleblower ein Insider Threat, der Confidentiality brach.
\end{itemize}
\block{PRISM, TEMPORA, XKEYSCORE}{F38--40,46--47; T47:22--51:34,58:58--1:00:32 \quad \textbf{?F81}}
\begin{itemize}
\item \textbf{PRISM} = Zugriff auf die Server von Providern wie Microsoft, Google, Yahoo, Facebook, YouTube, Skype, AOL und Apple, genannt \textbf{downstream collection} (\glqq Never mind the key, get the server\grqq).
\item \textbf{Warrant}: Für US citizens braucht es probable cause (FBI), Foreigners (alle anderen, NSA/CIA) haben laut Vorlesung diesen Schutz nicht.
\item \textbf{TEMPORA} = großflächiger Zugriff auf Optical Fibers an Dutzenden Orten (\textbf{upstream collection}), dazu Computer Network Exploitation (bulk und targeted Hacking).
\item \textbf{XKEYSCORE} = verteilte Search Engine über ca.\ 100 Repositories weltweit für Anfragen wie \glqq show me everybody in Afghanistan who's communicating in German\grqq.
\item \textbf{Routing}: Traffic nimmt den cheapest, nicht den physisch direktesten Pfad, deshalb fließt viel Kommunikation durch die USA.
\end{itemize}
\block{Five Eyes: zwei Fallbeispiele}{F42--44; T51:34--56:34}
\begin{itemize}
\item \textbf{Google (Leaked Slide)}: SSL lief nur bis zum Google Front End (GFE), zwischen den Data Centers floss der Traffic im Klartext, weshalb Google danach auch diesen Verkehr verschlüsselte.
\item \textbf{Operation Socialist (2010--13)}: GCHQ hackte Belgacom, um Telefone in Belgien (EU) im industriellen Maßstab abzuhören, also EU-Officials und Telefonverkehr nach Afrika.
\item \textbf{Vorgehen laut Vorlesung}: Malware bei Belgacom-Technikern, Lateral Movement bis zum GRX-Roaming-Proxy und dann Man-in-the-Middle auf den Mobilfunkverkehr.
\end{itemize}
\block{China: Hacks}{F48--52; T1:01:16--1:08:54 \quad \textbf{?F81 ?F83}}
\begin{itemize}
\item \textbf{Strategic Peer Competitor} der USA, dessen Hacking sich von smart people mit simple tools (2000er) zu systematischen Operationen im industriellen Maßstab entwickelte.
\item \textbf{Hacks}: Dalai Lama (2008), Google (2009), Office of Personnel Management (OPM, 2015) und Microsoft (2021).
\item \textbf{Dalai Lama (2008)}: Spear-Phishing eines Mönchs in Dharamsala, dann Malware auf dem Mailserver, die PDF-Anhängen einen JavaScript Buffer Overflow einfügte und 35 der 50 PCs der Exilregierung übernahm.
\item \textbf{Google (2009)}: China hackte genau die Einrichtung, die Google für den warranted access von NSA/FBI bauen musste, um herauszufinden, welche eigenen Agenten in den USA dem FBI bekannt waren.
\item \textbf{OPM (2015)}: China stahl die Datenbank der Security Clearances mit allen Antworten auf intrusive Fragen, die zur Erpressung taugen (\glqq database of ruin\grqq).
\end{itemize}
\block{China: Belt and Road und Platform Advantage}{F49,53; T1:09:10--1:11:47 \quad \textbf{?F81}}
\begin{itemize}
\item \textbf{Belt and Road Initiative}: China hilft Entwicklungsländern beim Ausbau von Infrastruktur inklusive Telekommunikation (z.\,B.\ Huawei), die dann von China kontrolliert wird.
\item \textbf{Ziel}: einen eigenen Platform Advantage wie den der Five Eyes aufbauen, vor allem in Ländern, die nicht mit den USA verbündet sind.
\item \textbf{Scale Access} über Belt and Road, Huawei, ZTE, TikTok und AI (z.AI; Alibaba Cloud + Qwen).
\item \textbf{Full-stack competition}: Chips, Offshoring der Fertigung von US-Firmen und eigene Tech-Majors.
\end{itemize}
\block{Russia, Iran, North Korea}{F54--57; T1:11:47--1:16:39 \quad \textbf{?F81 ?F83}}
\begin{itemize}
\item \textbf{Kein Platform Advantage} $\rightarrow$ diese Staaten setzen auf Spear-Phishing, Hacking und Information Warfare.
\item \textbf{Iran}: Die Uran-Anreicherungszentrifugen wurden 2008--09 von USA \& Israel gehackt (Stuxnet), worauf Iran sich an den Saudis rächte.
\item \textbf{North Korea -- Sony Hack (2014)}: Reaktion auf den Film \emph{The Interview}, bei der interne Daten und E-Mails veröffentlicht wurden, laut Vorlesung eine Machtdemonstration.
\item \textbf{North Korea sonst}: Hauptziel ist laut Vorlesung das Stehlen von Cryptocurrencies.
\item \textbf{Russia}: Cyber-Warfare in regionalen Konflikten (Ukraine) und Einflussnahme auf Wahlen (US 2016, Romania 2025).
\end{itemize}
\block{Stuxnet: wie es funktionierte}{F55--56; T1:13:08--1:15:12 \quad \textbf{?F81 ?F83}}
\begin{itemize}
\item \textbf{1.\ Infection, 2.\ Search}: Eintritt per USB-Stick in Windows-Rechner, getarnt durch ein digitales Zertifikat, dann Suche nach dem Siemens Industrial Control System der Zentrifugen.
\item \textbf{3.\ Update}: Ist der Rechner kein Ziel, tut Stuxnet nichts, sonst lädt der Worm neuere Versionen von sich nach.
\item \textbf{4.\ Compromise, 5.\ Control}: Zero-Day-Vulnerabilities übernehmen die PLCs (Logic Controllers), Stuxnet spioniert zuerst und lässt dann die Zentrifugen bis zum Versagen drehen.
\item \textbf{6.\ Deceive and destroy}: Falsches Feedback an die Kontrolleure verhindert, dass jemand den Schaden rechtzeitig bemerkt.
\item \textbf{Bedeutung}: erste Advanced Persistent Threat gegen cyber-physische Infrastruktur (Zäsur in der OT Security) mit vielen Infektionspfaden (USB, Contractors, S7 Project Files), die das Programm um etwa zwei Jahre zurückwarf.
\end{itemize}
\block{Russia: Ukraine, SolarWinds, US-Wahl 2016}{F54,58--60; T1:16:39--1:22:28 \quad \textbf{?F81}}
\begin{itemize}
\item \textbf{Ukraine Power Grid (2015)}: Start über ein präpariertes Word-Dokument, rund eine Viertelmillion Menschen stundenlang ohne Strom und gleichzeitig DoS auf die Callcenter -- der erste große Blackout durch Hacker.
\item \textbf{NotPetya (2017)} sollte ukrainische Firmen treffen und breitete sich bis nach Europa aus, weitere Folienbeispiele sind 2014 Vote-Counting System, 2022 WhisperGate Ransomware und Angriffe auf Military/Government HQ.
\item \textbf{SolarWinds (2020)}: Supply-Chain Attack, bei dem Schadcode in eine DLL (SolarWinds.Orion.Core.BusinessLayer.dll) legitimer Software kam und mit ihr an die Kunden verteilt wurde.
\item \textbf{SolarWinds-Backdoor}: prüft die Umgebung (Defense Evasion), sammelt Infos, verbindet sich mit Command-and-Control, exfiltriert Daten und erlaubt Hands-on-Keyboard-Angriffe wie Credential Theft und Lateral Movement.
\item \textbf{US-Wahl 2016}: Russische Intelligence Officers hackten E-Mail-Konten der Clinton-Kampagne und das DNC, veröffentlichten das Material über Online-Personas und WikiLeaks und nutzten Troll Farms pro Trump.
\end{itemize}
\block{Crooks: Cybercrime Ecosystem}{F62--64; T1:23:50--1:24:45 \quad \textbf{?F80 ?F82 ?F83}}
\begin{itemize}
\item \textbf{Größenordnung}: UK 2019 rund 1 Million Property Crimes, aber etwa 2,5 Millionen Online-Frauds und Scams, weil Kleinkriminalität ins Netz wanderte.
\item \textbf{Ecosystem} aus Bulk Attackers, Targeted Attackers, Tool Providers und Cashout Gangs, wobei Money Mules Geld über Konten schieben, bis es nicht mehr rückverfolgbar ist.
\item \textbf{Criminal Infrastructure}: unregulierte Cryptocurrency Exchanges und Botnets.
\item \textbf{Summary}: seit ca.\ 2005 über Underground Forums organisiert, dominiert Property Crime (keine Nähe nötig, skaliert leichter) und ist international und spezialisiert.
\item \textbf{Seit 2020}: Ransomware ist der große Wachstumsbereich.
\end{itemize}
\block{Ransomware: Entwicklung}{F65--66 \quad \textbf{?F82 ?F83}}
\begin{itemize}
\item \textbf{Ransomware} = Malware, die die persönlichen Daten des Opfers verschlüsselt, bis ein Lösegeld gezahlt wird.
\item \textbf{Early Days}: einfache PC-Verschlüsselung als Bulk Attack (petty crime) auf Privatnutzer, Lösegeld um 1 Bitcoin ($\approx$ 300 GBP).
\item \textbf{Current Landscape}: ausgefeilte large-scale Attacks auf große Firmen und verwundbare Institutionen, Lösegeld von Hunderttausenden bis Millionen.
\item \textbf{Cryptocurrencies} sind schwer rückverfolgbar und damit vom aufkommenden Zahlungsmittel zum Critical Enabler der Ransomware-Ökonomie geworden.
\end{itemize}
\block{The Swamp: Motiv Abuse}{F67--71; T9:08--9:17 \quad \textbf{?F80}}
\begin{itemize}
\item \textbf{Terrorism Recruitment} und \textbf{Child Sex Abuse Material (CSAM)}.
\item \textbf{Hate Campaigns}, z.\,B.\ \#GamerGate über 4chan, Reddit, YouTube und 8chan, teils auch Hacktivism.
\item \textbf{Intimate Relationship Abuse}.
\item \textbf{School/Workplace Bullying}, Folienbeispiel: Facebook Leaks (2021, Frances Haugen).
\end{itemize}
\block{Social Media und Content Moderation}{F72 \quad \textbf{?F82}}
\begin{itemize}
\item \textbf{Nutzer}: Laut einer Oxford/TUM-Studie (2025) wollen die meisten Menschen mehr Moderation (79\,\%).
\item \textbf{Governments} erhöhen den Druck zu zensieren bzw.\ zu moderieren, z.\,B.\ EU Digital Services Act und UK Online Harms Bill.
\item \textbf{Tech Companies} filterten schon viel zu hohen Kosten (Sex Abuse, Terror, Hate Speech, Nudity), als Backlash hörten manche wieder auf (Twitter/X, Facebook), und AI ist ein double-edged sword.
\item \textbf{Like/Share} führte zu performative Shaming, Social Media wirkt als Outrage Machine mit Folgen für die Mental Health.
\item \textbf{Solutions} sind teils technisch, teils politisch und teils eine Frage sich entwickelnder sozialer Normen.
\end{itemize}
\block{AI-Agents: Tool or Actor?}{F73--74 \quad \textbf{?F83}}
\begin{itemize}
\item \textbf{Kernfrage}: Ist ein AI-Agent Adversary oder Tool, also hat agentic AI Moral Agency oder ist sie nur eine Verlängerung menschlicher Absicht (Virtue Ethics, Moral Agency Theory)?
\item \textbf{Traditional Philosophy} sagt nein und behandelt Tools als value-neutral.
\item \textbf{Conditions for moral responsibility}: Freedom (Handeln ohne Zwang) und Epistemic Competence (Absichten und Folgen verstehen).
\item \textbf{Complicated Supply Chains} verschleiern die Accountability im Ecosystem und erzeugen potenziell distributed moral actions.
\item \textbf{Safety vs.\ Security}: Die Folie verweist auf diese Unterscheidung aus Systemsicherheit.
\end{itemize}
\block{Querschnittsfrage: Online Threats erkennen und bekämpfen}{aus F33,64--65,72,74--75 zusammengesetzt \quad \textbf{?F83}}
\begin{itemize}
\item \textbf{Kein Totalschutz}: Schutz gegen alle Threats zu vertretbaren Kosten ist unmöglich, also priorisieren.
\item \textbf{Verfolgung schwer}: Cryptocurrencies und ein internationales, spezialisiertes Ecosystem erschweren die Zuordnung.
\item \textbf{Verteilte Verantwortung}: AI-Supply-Chains verschleiern Accountability, und wir alle wirken als Amplifier.
\item \textbf{Lösungen} sind teils technisch, teils politisch und teils soziale Normen.
\end{itemize}

\section*{2\quad Hinweise der Lehrperson}
\block{Was kommt zum Minitest?}{T1:15--1:41; F76,84; M; E-Mail}
\begin{itemize}
\item \textbf{Transkript}: \glqq I really want to stress that you have to learn the definitions for the mini tests\grqq{} -- jede Definition der Einheit kann drankommen.
\item \textbf{E-Mail der Lehrperson}: Minitest am Freitag, ca.\ 5 Minuten, den Foliensatz \emph{inklusive} der Fragenfolien lernen.
\item \textbf{F84}: Fokus auf die Fragen am Ende des Foliensatzes, vor allem Definitionen daraus, aber das ganze Skriptum lernen; 2--3 Fragen, Multiple Choice und Fill in Text, diesmal \textbf{keine} Mathematical Computations.
\item \textbf{Moodle}: Die Fragen am Foliensatzende sind relevant, aber nicht vollständig (F76: \glqq not exhaustive\grqq), und die Folien enthalten auch Details, die in der Vorlesung nicht besprochen wurden.
\item \textbf{Abdeckung}: Jede Frage der Fragenfolien F77--F83 ist oben einem Block zugeordnet (Markierung \textbf{?F77} \dots{} \textbf{?F83} rechts im Blockkopf).
\end{itemize}
\block{Abgrenzung und ausdrückliche Merk-Hinweise}{T9:28,35:07--35:37,39:32,42:26--42:37,1:18:16; F3}
\begin{itemize}
\item \textbf{Nicht in diesem Kurs}: Attack Vectors, Tactics und Techniques (OWASP-Sicht), weil das die Definitionen sprengen würde.
\item \textbf{\glqq That's for you to know\grqq}: die Hierarchie Protection $\rightarrow$ Protection Profile $\rightarrow$ Security Target $\rightarrow$ Implementation.
\item \textbf{Hacktivism und Script Kiddies} werden nicht behandelt, Researchers und Tools kommen in der nächsten Einheit.
\item \textbf{Nächste Einheit}: Security Policies und Targets mit formalen Modellen, laut F3 prüft jeder Minitest nur die vorige Einheit.
\item \textbf{\glqq I want you to remember\grqq}: Neben dem Ukraine-Blackout gab es NotPetya, das nur ukrainische Firmen treffen sollte und sich nach Europa ausbreitete.
\end{itemize}
\block{Logistik am Freitag}{F2--3,84; P9; T0:00--1:09,41:18; M}
\begin{itemize}
\item \textbf{Termin}: Minitest in den ersten 5 Minuten der Vorlesung am 23.10.2026, also pünktlich sein (F3: \glqq be on time\grqq).
\item \textbf{Medium unklar}: F84 und P9 sagen Moodle Quiz, F3 sagt PEN + PAPER, in Moodle ist noch kein Quiz sichtbar, also (eigene Empfehlung) Gerät mit Moodle-Login und Stift mitnehmen.
\item \textbf{Wertung}: 12 Minitests à 5 Punkte, die besten 9 zählen (max.\ 45), bei mehr als 3 Fehlterminen oder unter 20 Punkten ersetzt das Final Exam (22.1., 11:45) die Minitests.
\item \textbf{Anwesenheit}: Die Liste geht in der Mitte oder am Ende der Stunde herum, höchstens dreimal fehlen, und wer stört, fliegt raus.
\item \textbf{Ältere Preamble (P9)} nannte noch 2--4 Fragen sowie True/False und Mathematical Computations, F84 ist die neuere Fassung.
\end{itemize}
\block{Poster-Thema wählen}{T0:28--0:52; P11--14; M; E-Mail}
\begin{itemize}
\item \textbf{Auftrag laut E-Mail und Vorlesung}: ein Thema wählen, in Moodle über \glqq Task: Select Topics for Written Report/Poster Session\grqq.
\item \textbf{First-come-first-served}: Die Plätze pro Thema sind begrenzt (meist 1), ein eigenes Thema per E-Mail mit der Lehrperson absprechen, die über die Eignung entscheidet.
\item \textbf{Report}: 8 Seiten zum gewählten Thema mit Overleaf-Template und 1--5 zitierbaren Quellen (kein Wikipedia), Abgabe bis 28.02.2027 (soft deadline mit Abzug).
\item \textbf{Poster Session} mit final minitest am 15.1.2027 um 11:45 (Moodle; F2 nennt 16.1.), Anwesenheit Pflicht, sonst 25-Minuten-Ersatzpräsentation in einem verwandten Kurs.
\item \textbf{Poster-Druck}: PDF (300 dpi, A2 oder A1) bis 9.1. schicken, dann druckt die Lehrperson über die Kopierzentrale.
\end{itemize}

\section*{3\quad Gut zu wissen}
\block{Details aus den Bildfolien}{F39--40,47}
\begin{itemize}
\item \textbf{PRISM-Zeitleiste}: Microsoft zuerst (9/11/07), Apple zuletzt (Oct 2012), Kosten ca.\ $\$$20M pro Jahr.
\item \textbf{PRISM-Daten}: E-Mail, Chat (Video, Voice), Videos, Photos, Stored Data, VoIP, File Transfers, Video Conferencing, Notifications of target activity und Social-Networking-Details.
\item \textbf{XKEYSCORE (2008)}: ca.\ 150 Sites mit über 700 Servern weltweit.
\end{itemize}
\block{Bewertung: Angaben widersprechen sich}{F3; P7--8; T0:00--0:28; M}
\begin{itemize}
\item \textbf{Punkte}: Aktuell (F3, Transkript) Minitests bzw.\ Final Exam 45, Report 25, Poster/Presentation 30, die ältere Moodle-Seite nennt noch 45/45/10.
\item \textbf{Bestehen}: Moodle verlangt mindestens 20 Punkte auf Minitests oder Final Exam (P8: 22) und insgesamt mindestens 50 Punkte.
\item \textbf{Notenschlüssel}: 0--49,9 = 5, 50--54,9 = 4, 55--74,9 = 3, 75--84,9 = 2, 85--100 = 1.
\end{itemize}

\section*{4\quad Selbsttest -- zehn Fragen im Stil von F77--F83}
\qa{1}{What is Security Engineering?}{Systeme bauen, die trotz \textbf{malice, error and mischance} dependable bleiben (design, implement, test, adapt); ohne Threat Actor wäre es Safety.}
\qa{2}{Describe the four key steps in Security Engineering Basics.}{\textbf{Threat Model $\rightarrow$ Security Policy $\rightarrow$ Implementation $\rightarrow$ Assurance}: Gegner bestimmen, Schutzziele festlegen, Mechanismen umsetzen, Erfolg auch Dritten beweisen.}
\qa{3}{Define subject, principal and person.}{\textbf{Person}: physical oder legal person; \textbf{Subject}: physical person in irgendeiner Rolle; \textbf{Principal}: authentifizierbare, im Security System handelnde Entität, auch Gerät, Role, Channel oder Compound.}
\qa{4}{What is the difference between trust and trustworthiness?}{\textbf{Trusted}: Versagen kann die Security Policy brechen (z.\,B.\ TPM); \textbf{trustworthy}: versagt nicht (Gegenbeispiel: Key Material verkaufender NSA-Mitarbeiter).}
\qa{5}{What is the difference between confidentiality and privacy?}{\textbf{Confidentiality}: Pflicht, fremde Geheimnisse zu schützen (Nutzen der Organisation); \textbf{Privacy}: Recht/Fähigkeit, eigene Daten und Raum zu schützen (Nutzen des Individuums).}
\qa{6}{Explain integrity and authenticity.}{\textbf{Integrity}: Daten wurden nicht verändert (at move/rest); \textbf{Authenticity}: \textbf{Integrity + Freshness + Proof of Identity}, also ein genuine Principal und kein Replay.}
\qa{7}{What is a weakness / vulnerability / exploit? CVE/CWE?}{\textbf{Weakness}: kann zu Vulnerabilities beitragen (CWE-125); \textbf{Vulnerability}: führt mit einer Threat zum Security Failure (CVE-2014-0160 Heartbleed); \textbf{Exploit}: Methode/Code zur Ausnutzung.}
\qa{8}{How does a security target differ from a security policy? What is a protection profile?}{\textbf{Policy} = what (Business Goals); \textbf{Target} = how (Spezifikation für ein Produkt); \textbf{Protection Profile} = device-independent Interface dazwischen für vergleichende Evaluations.}
\qa{9}{What is PRISM, and how does it differ from TEMPORA and XKEYSCORE?}{\textbf{PRISM}: Serverzugriff bei großen Providern (downstream collection, \glqq get the server\grqq); \textbf{TEMPORA}: Glasfaser-Abgriff (upstream collection); \textbf{XKEYSCORE}: verteilte Search Engine über die gesammelten Daten.}
\qa{10}{What is Stuxnet? How did it work?}{Worm von \textbf{USA \& Israel} gegen Irans Zentrifugen (2008--09): USB-Infektion, Suche nach Siemens-Steuerung, Zero-Days übernehmen die PLCs, Zentrifugen drehen bis zum Versagen, falsches Feedback tarnt es; erste APT gegen cyber-physische Infrastruktur.}
\end{document}
